\Needspace{10\baselineskip}
\SourceLeaf{92}%
\noindent man hat so folgende Gleichungen
\[
\begin{aligned}
B^{(0)}&=\frac{\frac{3}{2}}{2h-\frac{1}{2}B^{(1)}},\qquad
B^{(1)}=\frac{\frac{5}{4}}{2h-\frac{3}{4}B^{(2)}},\\[6pt]
B^{(2)}&=\frac{\frac{7}{6}}{2h-\frac{5}{6}B^{(3)}},\qquad
B^{(3)}=\frac{\frac{9}{8}}{2h-\frac{7}{8}B^{(4)}}\quad\text{u.\,s.\,w.}
\end{aligned}
\]

\Needspace{27\baselineskip}
Man bedient sich dieser Formeln, um die Grössen $B^{(0)}$, $B^{(1)}$,
$B^{(2)}$, $B^{(3)}$ u.\,s.\,w.\ in verkehrter Ordnung aus einer vom
Anfang entferntern, etwa der letzten, welche man nöthig hat, abzuleiten. Diese zu finden
dient die Formel

\[
B^{(k)}=\frac{2k+3}{2k+2}t\mathord{.}\cfrac{1}{1+\cfrac{\dfrac{3}{(2k+2)(2k+4)}\,tt}{1-\cfrac{\dfrac{(2k+1)(2k+5)}{(2k+4)(2k+6)}\,tt}{1-\cfrac{\dfrac{1.5}{(2k+6)(2k+8)}\,tt}{1-\cfrac{\dfrac{(2k+3)(2k+7)}{(2k+8)(2k+10)}\,tt}{1-\cfrac{\dfrac{3.7}{(2k+10)(2k+12)}\,tt}{1-\text{etc.}}}}}}},
\]

\Needspace{27\baselineskip}
\noindent also z.\,B.
\[
B^{(9)}=\tfrac{21}{20}t\mathord{.}\cfrac{1}{1+\cfrac{\dfrac{3}{440}\,tt}{1-\cfrac{\dfrac{437}{528}\,tt}{1-\cfrac{\dfrac{5}{624}\,tt}{1-\cfrac{\dfrac{75}{104}\,tt}{1-\cfrac{\dfrac{1}{40}\,tt}{1-\cfrac{\dfrac{207}{320}\,tt}{1-\cfrac{\dfrac{45}{1088}\,tt}{1-\text{etc.}}}}}}}}}
\]
\[
B^{(11)}=\tfrac{25}{24}t\mathord{.}\cfrac{1}{1+\cfrac{\dfrac{1}{208}\,tt}{1-\cfrac{\dfrac{621}{728}\,tt}{1-\cfrac{\dfrac{1}{168}\,tt}{1-\cfrac{\dfrac{145}{192}\,tt}{1-\cfrac{\dfrac{21}{1088}\,tt}{1-\cfrac{\dfrac{93}{136}\,tt}{1-\cfrac{\dfrac{5}{152}\,tt}{1-\text{etc.}}}}}}}}}.
\]
